% Generated draft: requires mathematical and source review.
\documentclass[11pt]{article}
\usepackage[margin=25mm]{geometry}
\usepackage{fontspec}
\setmainfont{Latin Modern Roman}
\usepackage{amsmath,amssymb,mathtools,unicode-math}
\setmathfont{Latin Modern Math}
\usepackage{enumitem,needspace,xurl,hyperref}
\setlength{\parindent}{0pt}
\setlength{\parskip}{5pt}
\setlength{\emergencystretch}{3em}
\newcommand{\N}{\mathbb{N}}
\newcommand{\Z}{\mathbb{Z}}
\newcommand{\Q}{\mathbb{Q}}
\newcommand{\R}{\mathbb{R}}
\newcommand{\C}{\mathbb{C}}
\newcommand{\F}{\mathbb{F}}
\newcommand{\A}{\mathbb{A}}
\newcommand{\PP}{\mathbb{P}}
\newcommand{\E}{\mathbb{E}}
\newcommand{\eps}{\varepsilon}
\newcommand{\dd}{\,\mathrm{d}}
\begin{document}
\section*{UnsolvedMath ID: 10600010}
\textbf{Source problem code:} catalogue entry
\subsection{Definitions and mathematical statement}
The compound item has three distinct parts.
\begin{enumerate}
\item For a virtual link \(L\), let \(B(L)\) denote its fundamental biquandle, presented in the usual manner by semiarc generators and the biquandle crossing relations. The problem asks whether \(B(L)\) classifies virtual links up to the source's mirror-image identifications: does \(B(L)\cong B(L')\) force \(L'\) to be virtually equivalent to \(L\), apart from the permitted mirror symmetries? The source explicitly records the known coincidence
\[B(L)\cong B(rm_\perp L),\]
where \(m_\perp\) is the mirror whose plane is perpendicular to the diagram plane and \(r\) reverses orientation. This specific simultaneous operation must be included; the excerpt does not identify every kind of mirror/reversal or state that each independently preserves the invariant. The biquandle convention and whether extra component, peripheral or orientation data are retained also require definition.
\item Find good examples of weak biquandles which are not strong. The source expressly says \emph{this problem is solved}, but its citation and the weak/strong definitions are elided. The historical request and supplied solved qualification are both retained; no new existence conjecture is asserted for this part.
\item Study the algebra associated with the linear case, having two generators \(A,B\) and the displayed relation
\[[B,(A-1)(A,B)]=0.\]
The ground ring, whether the generators are units, meaning of \((A,B)\), and exact algebraic information sought are absent. Under the usual additive-commutator reading, \([U,V]=UV-VU\), but the parentheses \((A,B)\) cannot silently be assigned that same meaning or a multiplicative commutator. Once these conventions are fixed, the task concerns the algebra presented by this relation, not an arbitrary commutative two-generator algebra.
\end{enumerate}
A virtual link is represented by diagrams with classical over/under crossings and virtual crossings, modulo the generalized Reidemeister/detour moves. The source's classification request concerns equivalence of these links, not just equivalence of their underlying unlabeled crossing graphs. Missing embedded references prevent a complete axiomatic definition of its particular fundamental, weak and strong biquandles.
\subsection{Short English statement}
Determine whether the fundamental biquandle completely classifies virtual links modulo the intended mirror symmetry, retain the already-solved weak-versus-strong example question, and study the displayed two-generator algebra with its unresolved notation.
\subsection{Sources}
\begin{itemize}
\item Local UnsolvedMath registry, permanent ID \(10600010\), exact source code \texttt{catalogue entry}.
\item \emph{Unsolved Problems in Virtual Knot Theory and Combinatorial Knot Theory}, revised survey, problem-list item 10: \url{https://arxiv.org/abs/1409.2823}. Registry access 2026-07-29; source-tex extraction. Embedded citations for the biquandle, its mirror property, the solved subproblem and the algebra are all elided in the supplied text.
\end{itemize}
\end{document}
